\documentclass[11pt]{beamer}
\usepackage[utf8]{inputenc}
\usepackage[T1]{fontenc}
\usepackage[french]{babel}
\usepackage{booktabs}
\usepackage{tikz}
\usetikzlibrary{arrows.meta,positioning}
\usetheme{Madrid}

%% pdflatex histoiresecretepresentation.tex && pdflatex histoiresecretepresentation.tex && rm *.log *.aux *.nav *.out *.snm *.toc

\newcommand{\carreau}{\ensuremath{\diamondsuit}}
\newcommand{\coeur}{\ensuremath{\heartsuit}}
\newcommand{\pique}{\ensuremath{\spadesuit}}
\newcommand{\trefle}{\ensuremath{\clubsuit}}

\title{L'Histoire secrète et le tarot}
\subtitle{Enseignes, correspondances et discordances}
\author{Synthèse documentée --- série de Jean-Pierre Pécau (Delcourt)}
\date{\today}

\begin{document}

% ---------------------------------------------------------------
\begin{frame}
\titlepage
\end{frame}

% ---------------------------------------------------------------
\begin{frame}{Le système des Archontes}
\begin{itemize}
\item Vers 3000 av.~n.~è., le chaman Père Loup remet à quatre petits-enfants
quatre \og ivoires \fg{} : un seul ivoire par archonte, jamais réunis.
\item Les quatre deviennent immortels, à la tête de leurs \og maisons \fg{}
(Coupes, Lances, Épées, Deniers), et manient la \og rétro-synchronicité \fg{}
(inspirée de la synchronicité de Jung) pour infléchir le hasard.
\item Chaque archonte possède un triple ancrage :
\begin{itemize}
\item une \textbf{enseigne latine / tarot} (l'ivoire, la maison) ;
\item une \textbf{couleur française moderne} (\carreau, \coeur, \pique, \trefle) ;
\item un \textbf{atout majeur} du tarot de Marseille.
\end{itemize}
\end{itemize}
\end{frame}

% ---------------------------------------------------------------
\begin{frame}{Correspondances officielles de la série}
\begin{center}\small
\begin{tabular}{lllll}
\toprule
\textbf{Archonte} & \textbf{Couleur} & \textbf{Enseigne} & \textbf{Atout} & \textbf{Maison} \\
\midrule
Dyo & \coeur{} Cœur & Coupe & Le Pape (V) & Coupes \\
Reka & \pique{} Pique & Bâton / Lance & La Papesse (II) & Lances \\
Aker & \carreau{} Carreau & Épée & L'Impératrice (III) & Épées \\
Erlin & \trefle{} Trèfle & Denier / Bouclier & L'Empereur (IV) & Deniers \\
\midrule
G.~de Lecce & (Joker) & --- & Le Mat & Moines rouges \\
\bottomrule
\end{tabular}
\end{center}
\vspace{0.4em}
\small Les quatre atouts originels = les \textbf{arcanes majeurs II à V}
du tarot de Marseille (les figures d'autorité).\\
Le cinquième archonte (tome 3) incarne le \textbf{Mat}, arcane non numéroté,
extérieur au pacte initial.
\end{frame}

% ---------------------------------------------------------------
\begin{frame}{La dérivation historique des enseignes}
\begin{center}\small
\begin{tabular}{llll}
\toprule
\textbf{Latine} & \textbf{Germanique} & \textbf{Française} & \textbf{Témoignage} \\
\midrule
Épée & Feuille & \pique{} Pique & \emph{spades} $<$ \emph{espadas} \\
Denier & Grelot & \carreau{} Carreau & rond $\rightarrow$ losange \\
Bâton & Gland & \trefle{} Trèfle & \emph{clubs} $<$ \emph{bastos} \\
Coupe & Cœur & \coeur{} Cœur & universel \\
\bottomrule
\end{tabular}
\end{center}
\vspace{0.4em}
Filiation réelle : latines $\rightarrow$ germaniques $\rightarrow$ françaises
(simplification graphique des cartiers, fin XV\ieme{} siècle).
\end{frame}

% ---------------------------------------------------------------
\begin{frame}{Comparaison : la BD contre la tradition}
\begin{center}\small
\begin{tabular}{llll}
\toprule
\textbf{Enseigne} & \textbf{Couleur BD} & \textbf{Couleur classique} & \textbf{Verdict} \\
\midrule
Coupes (Dyo) & \coeur{} Cœur & \coeur{} Cœur & \textcolor{green!55!black}{conforme} \\
Lances / Bâtons (Reka) & \pique{} Pique & \trefle{} Trèfle & \textcolor{red}{inversé} \\
Épées (Aker) & \carreau{} Carreau & \pique{} Pique & \textcolor{red}{inversé} \\
Deniers (Erlin) & \trefle{} Trèfle & \carreau{} Carreau & \textcolor{red}{inversé} \\
\bottomrule
\end{tabular}
\end{center}
\vspace{0.4em}
\begin{itemize}\small
\item Seule Coupes $\rightarrow$ Cœur respecte la tradition.
\item \textbf{Épées $\rightarrow$ Carreau} : l'inversion la plus hétérodoxe
(épées $\rightarrow$ piques est quasi universel).
\item \textbf{Lances $\rightarrow$ Pique} : jeu de mots interne
(une pique \emph{est} une arme d'hast).
\item \textbf{Deniers $\rightarrow$ Trèfle} : rejoint l'école Chevalier \& Gheerbrant.
\end{itemize}
\end{frame}

% ---------------------------------------------------------------
\begin{frame}{Schéma : filiation historique vs attribution BD}
\begin{center}
\begin{tikzpicture}[>={Stealth[length=2mm]},every node/.style={font=\scriptsize}]
\node[align=right] (coupe)  at (0, 0)    {Coupe (Dyo)};
\node[align=right] (epee)   at (0,-1.2)  {Épée (Aker)};
\node[align=right] (baton)  at (0,-2.4)  {Bâton / Lance (Reka)};
\node[align=right] (denier) at (0,-3.6)  {Denier (Erlin)};
\node (coeurL)   at (8, 0)    {\coeur{} Cœur};
\node (piqueL)   at (8,-1.2)  {\pique{} Pique};
\node (carreauL) at (8,-2.4)  {\carreau{} Carreau};
\node (trefleL)  at (8,-3.6)  {\trefle{} Trèfle};
\draw[->] (coupe.east)  -- node[above]{hist.\ et BD} (coeurL.west);
\draw[->] (epee.east)   -- (piqueL.west);
\draw[->] (baton.east)  -- (trefleL.west);
\draw[->] (denier.east)  -- (carreauL.west);
\draw[->,red,dashed] (epee.east)   to[bend right=18] (carreauL.west);
\draw[->,red,dashed] (baton.east)  to[bend right=18] (piqueL.west);
\draw[->,red,dashed] (denier.east) to[bend right=18] (trefleL.west);
\node[align=left] at (4,-4.5)
  {--- filiation historique \quad \textcolor{red}{- -} attribution BD};
\end{tikzpicture}
\end{center}
\end{frame}

% ---------------------------------------------------------------
\begin{frame}{La controverse occultiste : Deniers / Bâtons}
Deux écoles s'affrontent depuis le XIX\ieme{} siècle sur les deux enseignes
\og rondes \fg{} :
\begin{itemize}
\item \textbf{Papus} (et la majorité des auteurs) :\\
Deniers $\rightarrow$ Carreaux, \quad Bâtons $\rightarrow$ Trèfles ;
\item \textbf{Chevalier \& Gheerbrant} (\emph{Dictionnaire des symboles}) :\\
Bâtons $\rightarrow$ Carreaux, \quad Deniers $\rightarrow$ Trèfles.
\end{itemize}
\vspace{0.6em}
La série (Deniers $\rightarrow$ Trèfle) rejoint donc la seconde école.\\
Mais \textbf{aucune tradition identifiée} ne soutient
Épées $\rightarrow$ Carreau : sur ce point, la série crée sa propre
cartomancie.
\end{frame}

% ---------------------------------------------------------------
\begin{frame}{Vocabulaire : la chronologie des ivoires}
Tome 1 (\emph{Genèse}), discours de Père Loup ($-3000$) :
\begin{itemize}
\item Dyo : l'ivoire de \textbf{Coupe} ;
\item Reka : l'ivoire de \textbf{Lance} ;
\item Aker : l'ivoire d'\textbf{Épée} ;
\item Erlin : \og l'ivoire du \textbf{bouclier} pour défendre les trois autres \fg{}.
\end{itemize}
\vspace{0.4em}
Le vocabulaire glisse ensuite selon les époques :
\begin{itemize}
\item \textbf{Bouclier} (forme la plus ancienne) $\rightarrow$ \textbf{Denier}
(ères monétaires) ;
\item \textbf{Lance} (forme guerrière) $\leftrightarrow$ \textbf{Bâton}
(nom tarotique traditionnel de la maison).
\end{itemize}
Ces glissements servent l'ancrage historique ; l'inversion des couleurs,
elle, ne correspond à aucune réalité historique documentée.
\end{frame}

% ---------------------------------------------------------------
\begin{frame}{Le Mat, l'Excuse et le Joker}
\begin{itemize}
\item \textbf{Le Mat} : arcane non numéroté du tarot de Marseille
$\rightarrow$ devient \textbf{l'Excuse} dans le jeu de tarot
(jouable à tout moment, hors hiérarchie).
\item Son écho moderne : le \textbf{Joker} des jeux de 52 cartes,
carte \og générique \fg{} échappant aux règles communes.
\item \textbf{Guillaume de Lecce} (5\ieme{} archonte, dès le tome 3, modelé
sur Guillaume III de Sicile) : archonte extérieur au pacte des quatre
ivoires --- \og la cinquième puissance, celle qu'on ne doit pas nommer \fg{}.
Un choix fidèle à la symbolique de sa carte.
\item Né torturé à Palerme, il est la \og conséquence directe \fg{} de la
tentative d'Aker de créer le \og roi du monde \fg{} (Frédéric II).
\end{itemize}
\end{frame}

% ---------------------------------------------------------------
\begin{frame}{Hiérarchie des arcanes mineurs : une autre liberté}
\begin{itemize}
\item Dans l'univers de Pécau, les \og soldats \fg{} sont classés en force
\textbf{croissante} : 2, 3, \dots, 9, as, valet, cavalier, roi/reine ---
l'as vaut plus que le 9.
\item Dans le \textbf{jeu de tarot réel}, la valeur est \textbf{décroissante} :
le 21 est le plus fort des atouts ; pour les enseignes, c'est le roi qui
domine.
\item La BD emprunte au tarot de Marseille son imagerie, mais sa hiérarchie
est une construction interne, proche des jeux de cartes modernes.
\end{itemize}
\vspace{0.6em}
\small Anecdote : une ancienne version de l'article Wikipédia indiquait
\og Aker = épées / \emph{Pique} \fg{} (la correspondance classique !),
corrigée depuis en \og Carreau \fg{} --- preuve de la force de l'association
mentale épées $\rightarrow$ piques.
\end{frame}

% ---------------------------------------------------------------
\begin{frame}{Confirmé, interprétatif, non tranché}
\textbf{Confirmé (sources concordantes)}
\begin{itemize}
\item Triplets enseigne / couleur / atout des cinq archontes ;
\item \og Lances \fg{} au lieu de \og Bâtons \fg{} dès le tome 1 ;
\item \og Bouclier \fg{} chez le chaman, \og Denier \fg{} ensuite ;
\item 5\ieme{} maison (moines rouges), 6\ieme{} (complexe militaro-industriel) ;
\item Inversion des couleurs vs tradition.
\end{itemize}
\textbf{Interprétatif}
\begin{itemize}
\item Le caractère \og volontaire \fg{} des inversions ; l'idée des quatre
atouts comme \og ordre temporel et spirituel \fg{} face au chaos du Mat.
\end{itemize}
\textbf{Non trouvé}
\begin{itemize}
\item Aucun débat de fans (blogs, forums) sur l'inversion ;
vérification primaire : les albums eux-mêmes.
\end{itemize}
\end{frame}

% ---------------------------------------------------------------
\begin{frame}{Sources}
\small
\begin{itemize}
\item Wikipédia : \og L'Histoire secrète \fg{}, \og Arcane majeur \fg{},
\og Tarot Rider-Waite \fg{}, \og Pique \fg{}, \og Carreau \fg{},
\og Trèfle \fg{}, \og Excuse \fg{}, \og Joker \fg{}.
\item Miroir fr-academic (ancienne version de l'article --- \og Aker = Pique \fg{}).
\item Babelio : critiques du tome 1 (citation du discours de Père Loup).
\item Top Comics, BDThèque : fiches et avis concordants.
\item Blog \emph{Les Secrets du Tarot} : controverse Papus vs Chevalier \& Gheerbrant.
\item Gaudry.be : 5\ieme{} et 6\ieme{} maisons, univers des arcanes.
\end{itemize}
\end{frame}

\end{document}